This appendix records how the work on this volume changed the note. The Details were first written against the version of the note at commit \texttt{18a12f9}. Where an expansion did not go through as written, or found a compressed or imprecise step, the note was corrected; the Details now describe the corrected note. The table lists these corrections for the chapters on Parts~VII--IX and for the parameter ledger, with the Detail that prompted each one. No main theorem changed. Unless the table says otherwise, the item is on the main line.

{\small
\begin{longtable}{@{}>{\raggedright\arraybackslash}p{0.15\textwidth}>{\raggedright\arraybackslash}p{0.25\textwidth}>{\raggedright\arraybackslash}p{0.54\textwidth}@{}}
\toprule
Detail & place in the note & what changed\\
\midrule\endhead
\bottomrule\endfoot
\multicolumn{3}{@{}l}{\emph{Part~VII}}\\
\ref{d7:defpar}, \ref{d7:A6table} & Definition~\Nref{def:avg-parameter}(b), Lemma~\Nref{lem:cov-avg}(c), Remark~\Nref{rem:general-averaging}(iii), (A9) and (A11) of Lemma~\Nref{lem:uniformity} & GAP for general averaging pairs, off the main line. The definition of the parameter of a pair lacked the transposed sup bound on the defect $E$, which (A9) and (A11) use; it is not implied by the other bounds. Added, with its proof for the model pair. Ba\l aban's pair has $E=0$, so no later use was affected.\\
\ref{d7:A6table} & Lemma~\Nref{lem:uniformity}, (A11) and Step~2 of the proof; Remark~\Nref{rem:general-averaging}(i) & (A11) names the forward bounds of $\bQ$ and $E$ used for $\bC$; the reality of the pair at real $U$ became part of property~(i), and the proof says where it is used. The Detail first listed every source line of Parts~II and~III that contains the pair; that list is now the table of the Detail, by result and step.\\
\ref{d7:A7} & Step~3 of the proof of Lemma~\Nref{lem:uniformity} & Wording: Step~2 of the proof of Lemma~\Nref{lem:local-coercive} compares with the model pair directly, so that Step~3 there concerns the model pair only. The earlier wording suggested a comparison with the general pair at $U=1$. General pairs only, off the main line.\\
\ref{d7:sharp} & Proposition~\Nref{prop:sharp-curvature} & GAP (trivial): the proof set $n_\delta=4R+16$, which need not be $\ge n_0$, and used $a_\delta\le a_1$ without stating it. Now $n_\delta=\max(n_0,4R+16)$ and $a_\delta\in(0,a_1]$.\\
\ref{d7:setting} & the Setting of Section~\Nref{sec:coarse-action} & Wording: the note wrote $\mathcal S(B)=S(U_k(e^{\ii B}W))$ for all real $|B|_\infty<2r_1$, where $e^{\ii B}W$ need not lie in $\mathfrak U_k(a_1)$. Now $\mathcal S(B)$ is characterized as a minimum, and the identity is stated only for $e^{\ii B}W\in\mathfrak U_k(a_1)$; the uses lie in this range.\\
\ref{d7:image}, \ref{d7:gfinv} & before Lemma~\Nref{lem:gf-step-invariance} & The sense of the $\delta$-identities is stated: identities of measures, and, where $\mathfrak a'$ is a submersion, pointwise identities of the fibre densities.\\
\ref{d7:coarse} & Proposition~\Nref{prop:coarse-action}(a),(c); proof of Theorem~\Nref{thm:variational} & The condition $C_6c_6\le\min(c_3,c_5)$, used in Proposition~\Nref{prop:coarse-action}, was not imposed; it is now part of the choice of $c_6$ (see also Detail~\ref{det:ledger-c6c7}).\\
\ref{d7:lc} & Lemma~\Nref{lem:local-comparison} & The lemma was stated for $z\in\T$, while Theorem~\Nref{thm:step-coercive} applies it at centres that need not be lattice points. It is now stated for any point of the torus.\\
\ref{d7:coercive} & proof of Theorem~\Nref{thm:step-coercive} & The threshold for the decay of $\boldsymbol\Gamma_\pi-a$ is part of the choice of $\eps_L$.\\
\ref{d7:A6table} & Definition~\Nref{def:avg-parameter}(b); Step~2 of the proof of Lemma~\Nref{lem:uniformity} & Second integrated audit, Round~D (8 November 2026), VII-N01. A correction of this volume: item~(i) said that the transposed sup bound of the defect is not implied by Definition~\Nref{def:avg-parameter}(b), which states it. It now says: not implied by the range, the forward bound and the norm bound of $E$ in~(b). The note did not change.\\
\ref{d7:image}, \ref{d7:gfinv} & Section~\Nref{ss:31.1}; the paragraph before Lemma~\Nref{lem:gf-step-invariance} and its proof of~(iv) & Third integrated audit, Round~E (9 November 2026), R09. A correction of this volume: Detail~\ref{d7:image} took the fibre integral as the value of the density at every $V'$. For an integrable density this holds almost everywhere; it holds at every $V'$ for a bounded Borel density that vanishes outside a compact subset of the submersion region, as in the small-field step. The pointwise identities hold where the fibre integrals converge absolutely: Detail~\ref{d7:gfinv} runs the Fubini argument on the fibre under $\int|f|\,d\nu_{V'}<\infty$, and its test function vanishes outside the set of the $V'$ of Lemma~\Nref{lem:composition}. The note now says the same, and its proof of~(iv) is the measure-level argument of the Detail.\\
\ref{d7:step1} & Step~1 of the proof of Lemma~\Nref{lem:uniformity} & The November~10 audit, Round~F (10 November 2026), F-R03-006. The dependence of $c_g$ on $a$ is now stated exactly: through $a$ itself and the bound of Proposition~\Nref{prop:R}(e); $C_B$ and the separation constant are absolute. The note's Step~1 says the same.\\
\ref{d7:defpar}, \ref{d7:step0} & Definition~\Nref{def:avg-parameter}; Step~0 of the proof of Lemma~\Nref{lem:uniformity} & Round~F, F-R07-002 and F-R07-003. (B) on $\Lam_4(c_y)$ in Definition~\Nref{def:avg-parameter} now refers to the bonds and differences that involve only points of $\Lam_4(c_y)$; in this sense every gauge of Definition~\Nref{def:regular} on a cube containing $\Lam_4(c_y)$ satisfies (B) with parameter $\eps$ there (a statement edit, in both notes). Detail~\ref{d7:defpar} uses this convention. Detail~\ref{d7:step0} now treats the tori of side less than $8M_0$: the trivial partition, and one real gauge on $\Tf$ for every cube.\\
\ref{d7:step34} & Steps~3 and~4 of the proof of Lemma~\Nref{lem:uniformity}; Lemma~\Nref{lem:bootstrap-local} & Round~F, F-R05-006. Lemma~\Nref{lem:bootstrap-local} now takes small inputs that are invariant under real gauge transformations and bounds the full covariant gradient (a statement edit, in both notes); Theorems~\Nref{thm:reg-gauge}(b) and~\Nref{thm:reg-min} carry the full gradient too. The note and the Detail say how Step~3 uses the lemma.\\
\ref{d7:image}, \ref{d7:gfinv} & Section~\Nref{ss:31.1}; Lemma~\Nref{lem:gf-step-invariance} & Round~F, F-R10-001. A correction of this volume: the every-$V'$ case of Detail~\ref{d7:image} now gives its reason. The normal Jacobian is bounded below on the compact support, and finitely many projection charts bound the volume of each fibre there; for the small-field integrands, $|Z|_\infty\le w$, the Haar density is bounded and the Jacobian of the elimination is bounded away from zero. The pointwise invariance of Detail~\ref{d7:image} and the identities of Detail~\ref{d7:gfinv} are stated for the absolutely convergent fibre-integral values. The note was right (Round~E).\\
\ref{d7:gf} & the paragraph after Definition~\Nref{def:gauge-fixing} & Round~F, F-R10-002. A fixed gauge-fixing strength $a_{\rm gf}$ must be positive. It multiplies the whole exponent $\Phi_W$, its remainder $G_3$ and the quadratic term, and the normalization changes by a background-independent factor. The note and the Detail now say so. The construction uses strength~$1$.\\
\ref{d7:onestep} & the paragraph ``Other references'' after Lemma~\Nref{lem:one-step-expansion} & Round~F, F-R06-004. The extension to lattices of spacings $\delta$ and $\delta'=L\delta$ needs $a\delta'^2\le c_2$, which follows from $a\le c_2$ for $\delta'\le1$. The note and the Detail now state it; Part~VII uses $\delta'=1$.\\
\ref{d7:coarse} & proof of Proposition~\Nref{prop:coarse-action}(c) & Round~F, F-R10-007. A correction of this volume: the Detail warned against a fine-lattice kernel bound of $d^2A_*$, which the proof does not use. It now says that no estimate of $d^2A_*$ is needed, and that bounds uniform in $\eta$ for the second derivatives in the coarse inputs follow from Theorem~\Nref{thm:variational}(f). The note's sentence now says that no bound on the second or higher derivatives of $A_*$ on the fine lattice is used.\\
\ref{d7:hessian} & Proposition~\Nref{prop:hessian} & Round~F, F-R10-008. A correction of this volume: the test curve of Step~(a) stays in $\mathfrak U(\theta_0)$ because $U'\in\mathfrak U(B_3\eps_1)$ by Theorem~\Nref{thm:balaban}(b), and $B_3a_1\le c_6/C^{\rm units}_L\le\theta_0/81$; the $(M,\eps)$-regularity of $U'$ alone bounds its curvature only by $C_{\rm curv}\eps$. Step~(b) bounds $\omega_0=\bH^TJ$ directly, without a Cauchy estimate in $B$. No constant changed; the ledger lists $C_{\rm curv}$ in S0.\\
\ref{d7:bloch} & the remark after Lemma~\Nref{lem:poincare-step}; Appendix~\Nref{app:numerics} & Round~F, F-R10-005. The numerical minimum $0.030$ of the Bloch form and its reciprocal $33$ are taken over the sampled momenta; they certify neither the minimum over all momenta nor a uniform Poincar\'e constant. The note and the Detail now say so. No choice of a constant uses these values.\\
\ref{d7:lc} & proof of Lemma~\Nref{lem:local-comparison} & Round~F, F-R04-005. A correction of this volume: the difference $\boldsymbol\Gamma_\pi(U^\natural)-\boldsymbol\Gamma(1)$ decays at a quarter of a common rate $\gamma_0$, at most the rate of the difference of the two operators and the two chosen rates of their inverses, as the corrected Lemma~\Nref{lem:perturbed-inverse} gives; it had $\gamma/4$. Only the form norm is used.\\
\ref{d7:coercive} & proof of Theorem~\Nref{thm:step-coercive} & Round~F, F-R05-001. A correction of this volume: the double commutator of the IMS step has the kernel $+\sum_j(\chi_j(x)-\chi_j(x'))^2K(x,x')$; the sign was wrong. The bound is unchanged.\\
\ref{d7:covc}, \ref{d7:order} & Corollary~\Nref{cor:covariance}(c); the order of the constants of Part~VII & Round~F, F-R10-004. A correction of this volume: Step~3 took a Lipschitz constant on the disc of the final radius $c'_L$, which is then shrunk. Now a radius $r_*$ with $2r_*\le r_1/4$ and a kernel bound $A_L$ on the disc of radius $2r_*$ are fixed first; the Schwarz lemma gives the constant $K_L=2A_L/r_*$, and $c'_L\le\min(r_*,r_1/4,c_Lr_*/(16A_LS_{\gamma_1}))$. The complex rate of Step~4 uses $A_L$. Detail~\ref{d7:order} and the ledger (Detail~\ref{det:ledger-M0}, Table~C, S2) list $r_*$ and $A_L$ before $c'_L$. The note was right.\\
\ref{d7:step} & proof of Theorem~\Nref{thm:step}(f) & Round~F, F-R09-002. $U_k(e^{\ii\hat Z}W)$ and $e^{\ii\eta A_*(\hat Z)}U'$ differ by a real gauge transformation, which need not lie in $\mathfrak G_0$, and $\mathfrak E_k$ is invariant under all real gauge transformations. The note and the Detail now say so; they said ``the same orbit''.\\
\multicolumn{3}{@{}l}{\emph{Part~VIII}}\\
\ref{d8:A1} & proof of Proposition~\Nref{prop:dirichlet-hessian}(c); Lemma~\Nref{lem:coercive-general} & Compressed step: the bounds on $d^2\mathcal N(0)$ and on $W_2^X$ at complex perturbations were spread over the proofs of Lemma~\Nref{lem:hessian-gamma}, Proposition~\Nref{prop:complex-ext}(c) and step~(v) of the proof of Proposition~\Nref{prop:complex-ext}. The argument is now in the proof of Proposition~\Nref{prop:dirichlet-hessian}(c), with $\delta_*$ fixed once and $C_{\rm cmp}c''_L\le\delta_*$.\\
\ref{d8:range3}, \ref{d8:35a}, \ref{d8:35d} & Proposition~\Nref{prop:dirichlet-averaged}(a),(d) & The proof used the symmetry of $W$, which the class of range-$3$ potentials does not assume. It now uses $\Rea$-coercivity and $\bC^X[W]^T=\bC^X[W^T]$, so the statement holds as written.\\
\ref{d8:hessgamma} & proof of Lemma~\Nref{lem:hessian-gamma} & The citations for $\|d^2\mathcal N(0)\|\le C\eps$ and $|\omega_0|\le C\eps$ now name Lemma~\Nref{lem:div-curvature} and block Cauchy--Schwarz; the bound on $\omega_0$ no longer cites Lemma~\Nref{lem:sup-bounds}.\\
\ref{d8:lf-a} & Proposition~\Nref{prop:local-fluct}(a) & The seam sum loses half of the rate. The statement gives $C_Le^{-\gamma'(R_0-2L-10M)/2}$ and $R_0(L,M)\le C_LM$, with $R_0(L,M)$ explicit in the proof; the Schur constant holds also at the complex perturbations, and the factor $16$ gives $\Rea\Delta_Y\ge\frac1{16}c_L$ there.\\
\ref{d8:scalar}, \ref{d8:system} & Propositions~\Nref{prop:scalar-interior}, \Nref{prop:system-interior} & $\Psi$ and $A_v$ are the supremum of the pointwise sum, so the contraction factor $\frac12$ is exact.\\
\ref{d8:rescaled} & Lemma~\Nref{lem:rescaled} & The centre $x$ is a lattice point.\\
\ref{d8:cg-complex} & Lemma~\Nref{lem:coercive-general} & The four changes must have sum $C''_Lc''_L\le\frac18c_L$, with $C''_L$ absorbing $C_W$; the condition $C_{\rm cmp}c''_L\le\delta_*$ was added, and the statement names the conditions on $\eps$.\\
\ref{d8:localcov} & Theorem~\Nref{thm:local-covariance} & The rate in $m$ is $\gamma''/2$. Off the main line.\\
\ref{d8:qb} & proof of Lemma~\Nref{lem:quadratic-bound} & The H\"older quotient of the $\eta$-term needs the chain rule, now stated.\\
\ref{d8:covflat} & \Neqref{eq:cov-vs-flat} & Stated for $\sigma\le3$, which covers the use with $\sigma\le3\eps_a$ in Remark~\Nref{rem:bg-to-graded}.\\
\ref{d8:dc} & Remark~\Nref{rem:domain-constants} & $\eps_{\max}$ is defined by doubled real-reference constants, with an explicit list (see Detail~\ref{det:ledger-emax}).\\
\ref{d8:oneloop} & proof of Corollary~\Nref{cor:one-loop} & A free small constant $c_{\rm ol}$ replaces the coercivity constant $c_L$; the domain conditions $\eps_a\le\eps_{\max}$ and $C_{\rm cmp}\delta_a\le c''_L$ are stated, so $\eps_a,\delta_a$ depend on $L$ and $M$. Off the main line.\\
\ref{d8:interaction} & Theorem~\Nref{thm:interaction-local}, Proposition~\Nref{prop:map-compare} & $r_5$, $\gamma_5$ are defined and $\eps\le r_5$ is assumed; $r_4\le\theta_0$. Off the main line.\\
\ref{d8:35e} & proof of Theorem~\Nref{thm:reg-locality}(a) & Second integrated audit, Round~C (7 November 2026), III-M01. The proof said that the commutator $[\mathfrak N^X,\chi]$ of the non-local part with the cutoff has a decaying kernel; it contains $D^T[\mathfrak FS_X,\chi]$, which is not bounded uniformly in $\eta$. The proof now expands the commutator by the Leibniz rule over the terms $X_i^TM_iY_i$ and lets the propagators absorb the derivatives at the ends. A correction of this volume: the step ``(b) and (d) at $U_\delta$'' said that the commutators with $\chi$ change by bounded amounts; this holds for $[L,\chi]$ only, and the non-local part is now treated in the same factorized form.\\
\ref{d8:ix10} & Lemma~\Nref{lem:telescoping} & Round~D, VIII-M01. The locality clause of the lemma needs $\mathcal R(\Lam_{m-1})\subset\mathcal R(\Lam_m)$, which the Detail used and the lemma did not assume. The inclusion is now a hypothesis of the lemma (a statement edit, in both notes); every application satisfies it. Off the main line.\\
\ref{d8:range3}, \ref{d8:35d} & the Conventions of Section~\Nref{sec:dirichlet-averaged}; proofs of Propositions~\Nref{prop:dirichlet-averaged}(d) and~\Nref{prop:reg-complex} & Round~D, VIII-M02. The transpose $W^T$ of a range-$3$ potential was said to have the same properties as $W$, or to lie in the same class; its dependence on $U$ can reach $2r_W$. The note and the Details now list the properties that the bounds use (reality, covariance, support, norm, range in $A$, holomorphic extension), keep the halo through $\bC^X[W^T]=\bC^X[W]^T$, and bound the transposed block norms. The comparison of domains uses the range in $A$ only; the range in $U$ enters only the locality. No bound, halo or constant changed.\\
\ref{d8:A1} & proof of Proposition~\Nref{prop:dirichlet-hessian}(c) & Round~D, VIII-N01. A correction of this volume: Step~6 used $C_W\delta\le C(\eps+2\delta_*)$, which is false for the $\eps$-independent $C_W$. It now uses the sharp bound $C_1\mathrm s(\delta\!A)/\delta_*\le C_1$ of Step~5, as the note does.\\
Section~\ref{d8:imports}, \ref{d8:35e} & Lemma~\Nref{lem:perturbed-inverse}; proof of Proposition~\Nref{prop:dirichlet-averaged}(e) & Round~F, F-R04-005. The lemma gave the difference of two inverses a quarter of the smaller of the two inverse rates; this fails when the perturbation decays more slowly. The rate is now a quarter of $\gamma_0$, the minimum of the perturbation rate and two chosen rates of the inverses (a statement edit, in both notes). Import~I2 names $\gamma_0$, and Detail~\ref{d8:35e} reads its $\gamma$ as this common rate. No number changed.\\
\multicolumn{3}{@{}l}{\emph{Part~IX}}\\
\ref{d9:step-small} & Lemma~\Nref{lem:step-small} & The statement had $C_{\mathfrak N}C_G^2$ in $\theta$; the proof gives $C_{\mathfrak N}C_G$, and the square was not proved for $C_G<1$, which occurs at (L5). The statement is now linear in $C_G$, and the proof says why (``The power of $C_G$''). Either form gives the bound used in Corollary~\Nref{cor:weak-covariance}.\\
\ref{d9:local-ops} & Lemma~\Nref{lem:weak-local-ops}(a), Theorem~\Nref{thm:weak-inverse}(a) & The margin is $r+2r_0$; an earlier version stated $r+r_0$. Only margin bookkeeping was affected.\\
\ref{d9:fam-setting}, \ref{d9:cut-par} & Section~\Nref{sec:weakenable}, level~(L5) of Section~\Nref{sec:nested-walks}, Section~\Nref{sec:tiers} & The bound $R_v\ge R_0+2+2r_{\rm loc}$ was $2$ short for the sentence on the $2r_{\rm loc}$-neighbourhoods; now $R_0+4+2r_{\rm loc}$. The argument that the cubes $\Lam^v_j$ are rectangles no longer uses the hat. The uses were covered before.\\
\ref{d9:weak-bg} & proof of Theorem~\Nref{thm:weak-bg}(d) & The proof stated the class $\mathcal K$ with $B$ on $Z_i^{(r+1)}$, while the statement claims $Z_i^{(r)}$. The class is now stated with $Z_i^{(r)}$, with the reason: $B$ enters only through $\bH(s)B$.\\
\ref{d9:L4} & proof of Proposition~\Nref{prop:dirichlet-hessian}(c) & The bound on $W_2^{X_j}$ and its change at complex references, used at the repetition of (L4), were a compressed step; they are now in the proof (see Detail~\ref{d8:A1}).\\
\ref{d9:wi-a} & Section~\Nref{sec:nested-walks}, Theorem~\Nref{thm:weak-inverse} & The thresholds $M_1$ are fixed level by level, $M_1$ is the largest of them, and $C_v$ is among its arguments.\\
\ref{d9:nl-tiers} & Proposition~\Nref{prop:nl-tiers}, Theorem~\Nref{thm:tier-walk} & The rates at complex references are named (the fixed coercivity constants and the fixed rate of the differences of Proposition~\Nref{prop:complex-ext}), and the reason why $2\theta'$ bounds $\theta'$ there is written.\\
\ref{d9:one-loop} & proof of Proposition~\Nref{prop:one-loop-polymer}(b) & The condition for the background activities reads $\eps_a\le\eps_{\max}$ and $C_{\rm cmp}\delta_a<\delta_P$; the slip $\delta_a\le\delta_P$ was corrected.\\
\ref{d9:lp-scalar} & Lemma~\Nref{lem:rescaled-lp}, proof of Proposition~\Nref{prop:scalar-lp}(b) & The lemma is stated at lattice centres $x\in\eta\Z^d$, and the covering uses lattice centres.\\
\ref{d9:step-small} & proof of Lemma~\Nref{lem:step-small}, ``Counting'' & Yoh integrated audit, Round~B (7 November 2026), SH-WGT-001. The audit read the displayed $\theta$ as carrying $e^{\kappa\mathsf N_{\rm b}}$ on its first term only; the display has it on all three terms. The counting in the note now says that the non-local and seam terms carry $s^{X(j)}s^{n_\omega}$, of degree at most $\mathsf N_{\rm b}+\deg\omega$. A correction of this volume: the two bounds of Step~5 dropped $e^{\kappa|X(j)|}\le e^{\kappa\mathsf N_{\rm b}}$; they now carry it. No constant changed.\\
\ref{d9:wi-b}, \ref{d9:prod-compat} & Theorem~\Nref{thm:weak-inverse}, statement and proof of~(b); the paragraph after Definition~\Nref{def:compatible} & Round~B, SH-LOC-001. The proof of~(b) inserted $\one_{X'}+\one_{X\setminus X'}$, which is not the identity on gauge fields and on $\ell^2(\mathfrak B(X))$, and treated a difference of local cut terms, a differential operator, as a bounded factor on the seam. Now $Q=\one_X-\one_{X'}$, the telescope runs over bounded pieces, $\one_X$ on $2$-forms is defined, and Theorem~\Nref{thm:weak-inverse} assumes for~(b) the comparison of the local commutators $[\fA^X_{\rm loc},h_j]G^X_j$, with constants $(C_GC_{\rm loc},\gamma_0)$ (a statement edit). Corrections of this volume: Detail~\ref{d9:wi-b} had $Q=\one_{X\setminus X'}$, the identity $\one_{X'}\one_{X_j}=\one_{X'_j}$ on gauge fields, a split of the local commutator into pieces not bounded uniformly in $\eta$, a bare $D^T$ in a cross term, and, in Step~5, the term $[D,P]A_0P$ and the derivatives of the compressed $\Delta_G$. It now compares already differentiated factors, and covers the split of the (L3) commutator of Detail~\ref{d9:wstep}. The claim of Detail~\ref{d9:prod-compat} listed $D^T_X$ among the bounded local factors; it now admits only bounded ones. The rate $\gamma_0/32$, the seam width and the loss $2\delta$ are unchanged.\\
\ref{d9:L1}--\ref{d9:L5} & the levels of Section~\Nref{ss:41.3} & Round~B: the new hypothesis of Theorem~\Nref{thm:weak-inverse}(b) is proved at (L1), (L3) (both passes) and (L5), in the note and in the Details; it is trivial at (L2), (L4). Detail~\ref{d9:L3}, Step~2, no longer lists $D$, $D^T_X$ among the bounded local factors, and the introduction of Section~\ref{d9:levels} names the new check.\\
\ref{d9:lift}, \ref{d9:N-tier} & proofs of Lemmas~\Nref{lem:lift} and~\Nref{lem:N-tier} & Round~B: the item ``a difference of (local) cut terms'' was removed from the seams of Lemma~\Nref{lem:lift} and of Detail~\ref{d9:lift}; at (L1) the compressed difference of the local parts vanishes. The cross term of Lemma~\Nref{lem:N-tier} uses $\one_{\Tf}-\one_{X'}$, the complement on every space. Detail~\ref{d9:N-tier}, Step~(N3), no longer lists $D^T_Z$ among the bounded local operators.\\
\ref{d9:wi-a} & proof of Theorem~\Nref{thm:weak-inverse}(a) & Round~E, FA-R11-001. The proof in the note inferred that the right inverse is the inverse from invertibility results of particular applications; it now uses that $\ell^2(X)$ is finite-dimensional, as the Detail does. The Detail no longer refers to the coercivity theorems that the note cited.\\
\ref{d9:nl-tiers} & Proposition~\Nref{prop:nl-tiers}; step~(vii) of the proof of Proposition~\Nref{prop:complex-ext} & Round~F, F-R04-005. The fixed rate of the differences at complex references is a quarter of the minimum of the perturbation rate and the two chosen inverse rates, in the note and in the Detail. No number changed.\\
\multicolumn{3}{@{}l}{\emph{The parameter ledger}}\\
\ref{det:ledger-emax} & Remark~\Nref{rem:domain-constants}; proof of Proposition~\Nref{prop:complex-ext}, step~(vii) and ``Smallness'' & The note defined $\eps_{\max}$ through ``the worst case of [the] constants over complex references of size $\delta_c\le\bar\delta$''. That worst case is available only through $\delta_P$, which is chosen after $\eps_{\max}$. Now every threshold is computed with doubled real-reference constants, at fixed coercivity constants and rates, and Proposition~\Nref{prop:complex-ext} shows that these bound the true thresholds.\\
\ref{det:ledger-eps-M} & proof of Proposition~\Nref{prop:first-step}(b),(c) & Compressed step: the first-step thresholds at non-critical references were met only through the existential $\eps_a$, $\delta_a$ of Proposition~\Nref{prop:one-loop-polymer} at $k=0$ and the rule, already stated, that the flow uses the worse of the values of a constant at $k=0$ and $k\ge1$. The proof now states the choice of $\eps_a$, $\delta_a$ at $k=0$.\\
\ref{det:ledger-emax-list} & Remark~\Nref{rem:domain-constants}; proof of Proposition~\Nref{prop:one-loop-polymer}(b) & The list of $\eps_{\max}$ named ``those of Parts~IV and~VII that they use''; it is now explicit, with $\eps_{\rm B}$, the smallness of Corollary~\Nref{cor:parts-apply} and $\eps_*$. The slip $\delta_a\le\delta_P$ was corrected.\\
\ref{det:ledger-c6c7} & proof of Theorem~\Nref{thm:variational} & The condition $c_6\le\min(c_3,c_5)/C_6$, used in Proposition~\Nref{prop:coarse-action}, was not imposed; it was added. No cycle arises.\\
\ref{det:ledger-a1B3} & proof of Theorem~\Nref{thm:balaban} & The condition $B_3a_1\le c'_3$, used in Part~XI (proof of Theorem~\Nref{thm:ol-chart}), was added.\\
\ref{det:ledger-M0} & Section~\Nref{sec:nested-walks}, Remark~\Nref{rem:parameters}(4), Proposition~\Nref{prop:local-fluct} & $M_1$ is read level by level, with $C_v$ among its arguments; $R_0(L,M)\le C_LM$ is stated in Proposition~\Nref{prop:local-fluct}.\\
\ref{det:ledger-gamma} & Proposition~\Nref{prop:flow-F}, (F3) & A correction of this volume: the Detail kept the former name $C_P$ of the constant of the pure sources in $C_3$, read as a generic constant of the step; it now uses the note's name $C_{\rm pure}$. The change of the note is listed under Detail~\ref{d11b:F3} (Subsection~\ref{d98b}).\\
\ref{det:ledger-off} & Proposition~\Nref{prop:map-compare}, Theorem~\Nref{thm:interaction-local}, Corollary~\Nref{cor:one-loop} & Three GAPs off the main line: the contraction in Proposition~\Nref{prop:map-compare}(a) needed a condition on $\theta_0$ (added to the proof of Theorem~\Nref{thm:variational}, with $r_4\le\theta_0$); $r_5$, $\gamma_5$ were not defined; the proof of Corollary~\Nref{cor:one-loop} used $c_L$ where a free small constant is needed. All three are repaired. No main-line statement was affected.\\
Table~C of the ledger & Lemma~\Nref{lem:coercive-general}; proof of Proposition~\Nref{prop:dirichlet-hessian}(c) & A correction of this volume: the row of $c''_L$ placed it among the constants of $L$ alone (stage S2). It depends on $M$ through $\delta_*$ and $C_W$, and is fixed in S4 after $M$; rows for $\delta_*$ and $C_W$ were added. The order of choice has no cycle either way.\\
\ref{det:ledger-eps-M} & Lemma~\Nref{lem:local-comparison}; proof of Proposition~\Nref{prop:first-step}(c) & A correction of this volume (Yoh integrated audit, Round~A, 7 November 2026): condition~(1) of Step~1 had the radius $40M'$; the producer uses $40M'+1$ around a lattice point, and the Detail now does too. Step~3 read $\delta^{(0)}_a<\delta_{P,0}$; it now reads $C_{\rm cmp}\delta^{(0)}_a\le\delta_{P,0}$, as Detail~\ref{xb:d-k0-thresholds} has it. The note was right.\\
\ref{det:ledger-emax-list}, Table~D & proof of Proposition~\Nref{prop:one-loop-polymer}(b) & A correction of this volume (Round~A): Step~3 and the row of $\eps_a,\delta_a$ in Table~D allowed every $\delta_a<\delta_P$ and gave the example $\delta_a=\frac12\min(\delta_P,\delta_{P,0})$, which violates $C_{\rm cmp}\delta_a<\delta_P$ when $\delta_{P,0}=\delta_P$. They now require $4\delta_a<\delta_P$, and at $k=0$ also $4\delta_a<\delta_{P,0}$, with the example $\delta_a=\frac18\min(\delta_P,\delta_{P,0})$; Step~3 adds $c_a\le\delta_a/(4C_{\rm ext})$. The note was right.\\
\ref{det:ledger-torus} & Lemma~\Nref{lem:local-comparison}; Proposition~\Nref{prop:first-step}(b) & A correction of this volume (Round~A): the boxes of Lemma~\Nref{lem:large-cubes} have side at most $120M'+3$, not $120M'$.\\
\ref{det:ledger-c6c7}, \ref{det:ledger-a1B3} & proofs of Theorems~\Nref{thm:variational} and~\Nref{thm:balaban} & Round~C. A correction of this volume (V-W02): the status line of Detail~\ref{det:ledger-c6c7} attributed $c_7\le\eps_V$ to Section~\Nref{ss:25.2}; it is the ledger's reading of the thresholds that the proof does not name. In the note (``bounded by twice the supremum'') and in Detail~\ref{det:ledger-a1B3} (``the bound by twice the supremum'') (VI-N01), the text now reads $|\partial A|\le2|A|$ and $[\partial A]_\alpha\le2|\partial A|\le4|A|$. No constant changed.\\
Table~D & Proposition~\Nref{prop:complex-ext}; Section~\Nref{sec:step-quantities} & Round~E, FA-R13-002. A correction of this volume: the row of $C_P$ (step) said that $C_P$ depends on $\delta_P$ through $R_P$. $R_P$ is fixed in S4 and does not depend on $\delta_P$; the row now says so.\\
\ref{det:ledger-emax} & Remark~\Nref{rem:domain-constants}; Proposition~\Nref{prop:complex-ext} & Round~F, F-R04-005. The fixed rate is a quarter of the minimum of the perturbation rate and the two chosen inverse rates.\\
\ref{det:ledger-M0} & Remark~\Nref{rem:step-M}; Lemma~\Nref{lem:local-comparison} & Round~F, F-R10-011 and F-R10-004. A correction of this volume: the step on Part~VII gave the seam distance as $10M<9M'$, which fails for $M'<\frac{10}9M$. The distance of the seam from the inner cube is at least $20M'-M'-10M\ge9M'$, because $M'\ge M$; no ratio of $M'$ to $M$ is needed. The same step now chooses $c'_L$ after $r_*$ and $A_L$, and the complex part of $\gamma_C$ uses $A_L$.\\
Tables~A and~C; Subsection~\ref{subsec:ledger-order} & Theorem~\Nref{thm:chart}; Corollary~\Nref{cor:covariance} & Round~F, F-R07-007, F-R10-004 and F-R10-008. A correction of this volume: the row of $c_4$, $r_u$, $B_1$ now lists $a$, $\alpha$ and $M_0$, and not $M$; the row of $c'_L$ cites the bound of Detail~\ref{d7:covc}; the order lists $C_{\rm curv}$ in S0, and $r_*$, $A_L$, $c'_L$, then $\gamma_C$, in S2.\\
Former names & Subsection~\ref{subsec:ledger-names} & Not a correction: the note renamed the symbols listed there (all rows but the last, which renames conditions of this volume); this volume uses the new names throughout.\\
\end{longtable}
}
